% ============================================================================
% SECTION 2.6 : FORMALISATION DU TRILEMME CRO DANS LEAN4
% Source : partie2/cro_lean4.tex
% Version : 2.2 — alignée sur la formalisation Lean4 (section~\ref{sec:lean4})
% ============================================================================

\section{FORMALISATION DU TRILEMME CRO DANS LEAN4}
\label{sec:lean4}
\begin{sloppypar}
La section~\ref{sec:sotp_construction} a pr\'esent\'e les constructions de r\'ef\'erence
SOTP.
Il reste cependant \`a savoir si les preuves th\'eoriques \'etablies dans les
chapitres pr\'ec\'edents sont correctes. Cette question
\textbf{Obj~\ref{obj:obj5_verification}} motive l'\'etape ultime de notre progression
m\'ethodologique.

Nous v\'erifions formellement le Trilemme CRO conditionnel dans
l'assistant de preuve Lean4~\cite{Lean4,moura2015lean,Avigad2019}, constituant
l'\'etape \textbf{C (Certifier)} du processus F$\to$B$\to$E$\to$S$\to$\textbf{C}$\to$X.
Notre formalisation \'etablit la borne
$\gamma^{-1}(\Delta/\lambda + \kappa(\Delta+1)/\lambda^2 + \kappa^2/\lambda^3)$
et sa forme raffin\'ee. Elle fournit aussi deux t\'emoins de coh\'erence et de port\'ee, certifie
les cons\'equences d'un verdict d\'eterministe et constitue une biblioth\`eque Lean4
r\'eutilisable pour la cryptographie institutionnelle. Cette contribution \textbf{C5}
a aussi conduit \`a des formulations plus pr\'ecises de plusieurs \'etapes des
preuves publi\'ees, pr\'esent\'ees \`a la sous-section~\ref{subsec:lean4_erreurs}.

Les sous-sections suivantes pr\'ecisent d'abord ce que cette formalisation
couvre et ses limitations (section~\ref{subsec:portee_limitations}), puis les
conventions typographiques propres \`a cette section
(section~\ref{subsec:conventions_typographiques}). Elles rappellent ensuite les
\'el\'ements du Trilemme CRO s\'electionn\'es pour la formalisation, avec les points
critiques identifi\'es lors de l'analyse pr\'eliminaire
(section~\ref{subsec:elements_formaliser}), avant d'exposer les d\'efis techniques
de la formalisation (section~\ref{subsec:defis_formalisation}), le mod\`ele formel
et l'axiomatisation (section~\ref{subsec:modele_formel_lean4}) ainsi que la preuve
du c\oe ur alg\'ebrique (section~\ref{subsec:preuve_coeur_algebrique}). Les deux
derni\`eres \'enoncent le th\'eor\`eme principal v\'erifi\'e
(section~\ref{subsec:theoreme_principal_lean4}) puis appliquent la formalisation
aux familles de protocoles canoniques
(section~\ref{subsec:application_familles_canoniques}).

Le code source v\'erifi\'e est disponible en acc\`es ouvert (licence MIT,
\url{https://github.com/MINKAThierry/PhDThesis}~\cite{CROTrilemmaDepot}), avec une
version de Lean et de Mathlib fig\'ee (v4.34.0) et une int\'egration continue qui le
recompile \`a chaque modification.
\end{sloppypar}

\begin{boiteArbitre}
L'arbitre croyait que la preuve du Trilemme \'etait compl\`ete. La machine a
montr\'e que plusieurs \'etapes des preuves publi\'ees admettaient des formulations
plus pr\'ecises, retenues dans les versions pr\'esent\'ees ici. L'arbitre comprend
que la rigueur, comme l'opposabilit\'e, se v\'erifie, gr\^ace \`a la machine et \`a
la relecture qui se compl\`etent.
\end{boiteArbitre}

% -----------------------------------------------------------------------
\subsection{Port\'ee et limitations de la formalisation}
\label{subsec:portee_limitations}
% -----------------------------------------------------------------------
\begin{sloppypar}
Cette sous-section pr\'ecise ce que couvre la formalisation Lean4 pr\'esent\'ee dans
ce chapitre, ainsi que ses limitations explicites.

\subsubsection{Ce que la formalisation couvre}
\label{subsubsec:ce_que_couvre}

Notre formalisation couvre le c\oe ur alg\'ebrique du Trilemme CRO conditionnel.
Nous d\'efinissons une structure \texttt{ProtocolFamily} index\'ee par le
param\`etre de s\'ecurit\'e $\lambda$, avec des param\`etres institutionnels
fixes $(\Delta, \alpha, \beta, H(C))$. Nous
int\'egrons comme propri\'et\'es des instances de protocole les quatre pr\'emisses
du th\'eor\`eme, \`a savoir l'in\'egalit\'e de
Fano~\cite{Fano61,CoverThomas06}, l'in\'egalit\'e de Pinsker (sous-sous-section~\ref{subsubsec:pinsker_classique}) combin\'ee \`a la
r\'egularit\'e~\cite{Pinsker1964,Tsybakov2009}, l'hypoth\`ese d'extraction
institutionnelle~\cite{Simon82,Miller56} et la borne fiabilit\'e-information.
Nous prouvons formellement que, pour tout $I \le \kappa/\lambda$ (r\'egime de
pr\'eservation de la confidentialit\'e, Hypoth\`ese~\ref{ass:regime}), on a
\begin{equation}
\Conf \cdot \Rel \cdot \Opp \le \frac{1}{\gamma}\left(\frac{\Delta}{\lambda}
+ \frac{\kappa(\Delta+1)}{\lambda^2} + \frac{\kappa^2}{\lambda^3}\right)
\le \frac{\Delta}{\lambda}
+ \frac{\kappa(\Delta+1)}{\lambda^2} + \frac{\kappa^2}{\lambda^3},
\end{equation}
ainsi que la forme raffin\'ee~\eqref{eq:trilemme_cro_raffine}. Ces preuves
n'utilisent que l'alg\`ebre \'el\'ementaire et sont enti\`erement v\'erifi\'ees par
Lean4. Nous prouvons \'egalement que $(1/2)^\lambda$ est n\'egligeable, par un
passage \`a la limite (\texttt{Tendsto}) et un argument de filtre, ce qui comble
une lacune fr\'equente dans les expositions informelles.

\subsubsection{Ce que la formalisation ne couvre pas}
\label{subsubsec:ce_que_ne_couvre_pas}

Nous explicitons ci-dessous ce qui
n'est pas inclus dans cette version v\'erifi\'ee.

\textbf{Les axiomes (Fano, Pinsker) ne sont pas prouv\'es
ab initio.} Ils sont int\'egr\'es comme
propri\'et\'es des instances de protocole, laissant leur d\'erivation formelle \`a partir de la th\'eorie des
probabilit\'es comme travail futur.

\textbf{Les constructions SOTP ne sont pas encore couvertes par Lean4.}
La formalisation couvre le noyau alg\'ebrique du Trilemme, mais pas les
r\'eductions cryptographiques des constructions SOTP. Ces extensions
constituent une direction de recherche future.

\textbf{L'extension quantique du Trilemme CRO n'est pas couverte.}
La formalisation pr\'esent\'ee ici se limite au cadre classique. L'extension
aux syst\`emes post-quantiques n\'ecessite une
infrastructure th\'eorique distincte incluant les in\'egalit\'es de Pinsker
quantique, le lemme de mesure douce (sous-sous-section~\ref{subsubsec:gentle_measurement}) et les bornes de Fano
quantique. Cette extension est rel\'egu\'ee aux travaux futurs.
\end{sloppypar}

% -----------------------------------------------------------------------
\subsection{Conventions typographiques et linguistiques}
\label{subsec:conventions_typographiques}
% -----------------------------------------------------------------------
\begin{sloppypar}
Cette sous-section pr\'esente une sp\'ecificit\'e li\'ee \`a la nature de son
contenu. Le corps du texte, les d\'efinitions math\'ematiques, les
th\'eor\`emes et les explications utilisent la ponctuation et les conventions
typographiques fran\c{c}aises (espaces avant les deux-points, guillemets
fran\c{c}ais, etc.), en coh\'erence avec les autres chapitres de cette th\`ese.

En revanche, les listings de code Lean4 pr\'eservent la syntaxe anglaise de
ce langage de programmation pour garantir la validit\'e technique du code.
Les commentaires inclus dans ces listings sont r\'edig\'es en fran\c{c}ais pour
faciliter la compr\'ehension tout en maintenant la compatibilit\'e syntaxique.
Cette convention permet au code d'\^etre directement compilable par
l'assistant de preuve Lean4 tout en restant accessible pour un lecteur
francophone.
\end{sloppypar}

% -----------------------------------------------------------------------
\subsection{\'El\'ements du Trilemme CRO \`a formaliser}
\label{subsec:elements_formaliser}
% -----------------------------------------------------------------------
\begin{sloppypar}
Cette sous-section rappelle les \'el\'ements du Trilemme CRO conditionnel,
\'etabli \`a la section~\ref{sec:cro_classique}, qui seront formalis\'es dans
Lean4. Nous pr\'ecisons les simplifications op\'er\'ees par rapport \`a la
version th\'eorique, justifions le p\'erim\`etre retenu et identifions les
points critiques que la formalisation r\'ev\`ele.

\subsubsection{M\'etriques CLO \`a formaliser}
\label{subsubsec:metriques_formaliser}

Le cadre CLO d\'efinit trois
m\'etriques pour un protocole de preuve $\mathcal{P}$ avec espace de messages
$\mathcal{M}_\lambda$ ($|\mathcal{M}_\lambda| = 2^\lambda$), espace de
contexte $\mathcal{C}_\lambda$ et espace de preuves $\Sigma_\lambda$.

\begin{definition}[Confidentialit\'e]\label{def:conf_lean4}
\begin{equation}
\Conf(\mathcal{P}) = 1 - \max_{\substack{m_0,m_1\in\mathcal{M}_\lambda\\c\in\mathcal{C}_\lambda}}
   \mathrm{TV}(P_{\Sigma|M=m_0,C=c}, P_{\Sigma|M=m_1,C=c})
\end{equation}
o\`u $\mathrm{TV}$ est la distance de variation totale.
\end{definition}

\begin{definition}[Fiabilit\'e]\label{def:rel_lean4}
Soit $\hat{M}$ l'estimateur MAP de $M$ \`a partir de $(\Sigma, C)$. La fiabilit\'e
est alors d\'efinie par
\begin{equation}
\Rel(\mathcal{P}) = 1 - \Pr[\hat{M} \neq M]
\end{equation}
\end{definition}

\begin{definition}[Opposabilit\'e]\label{def:opp_lean4}
Soit $V: \Sigma \times \mathcal{C} \to \mathcal{Y}$ une fonction de
verdict. L'opposabilit\'e est alors d\'efinie par
\begin{equation}
\Opp(\mathcal{P}) = \frac{I(V;\Sigma|C)}{H(V|C)}
\end{equation}
o\`u $I$ et $H$ d\'esignent respectivement l'information mutuelle
conditionnelle et l'entropie conditionnelle.
\end{definition}

\subsubsection{Hypoth\`eses institutionnelles s\'electionn\'ees}
\label{subsubsec:hypotheses_formaliser}

Le trilemme repose sur quatre pr\'emisses, dont trois portent sur
l'institution interpr\`ete et le protocole, la quatri\`eme reliant fiabilit\'e et
information. Nous les formalisons telles quelles, sans simplification.
% Note de source : les labels LaTeX de ces quatre enonces ne correspondent pas
% a leur contenu ; seuls les titres font foi.

\begin{assumption}[Extraction born\'ee]\label{ass:fano_axiom}
Pour toute institution avec fonction de verdict $V$, on a
\begin{equation}
I(V;\Sigma|C) \le I(M;\Sigma|C) + \Delta
\end{equation}
o\`u $\Delta \ge 0$ quantifie la connaissance pr\'ealable ou la capacit\'e de
raisonnement au-del\`a du contexte formel $C$.
\end{assumption}

\begin{assumption}[Espace de verdict born\'e]\label{ass:pinsker_axiom}
La taille de l'espace de verdict satisfait
\begin{equation}
|\mathcal{Y}_\lambda| \le 2^{\alpha H(C)}
\end{equation}
pour une constante $\alpha > 0$ mesurant la granularit\'e de d\'ecision.
\end{assumption}

\begin{assumption}[R\'egularit\'e]\label{ass:extraction_axiom}
Un protocole $\mathcal{P}$ est $\beta$-r\'egulier si pour tout contexte $c$
et messages $m_0, m_1$, on a
\begin{equation}
D_{KL}(P_{\Sigma|m_0,c} \| P_{\Sigma|m_1,c}) \le \beta \cdot I(M;\Sigma|C)
\end{equation}
avec $\beta \ge 1$.
\end{assumption}

\begin{assumption}[Fiabilit\'e-information]\label{ass:reliability_axiom}
Pour toute instance de protocole, sans restriction sur $\Rel$, on a
\begin{equation}
\Rel \le \frac{I(M;\Sigma|C) + 1}{\lambda}
\end{equation}
La formulation initiale imposait la restriction $\Rel \ge 1/2$. La
formalisation a montr\'e qu'elle peut \^etre lev\'ee, car la majoration
$h(p_e) \le 1$ de l'entropie binaire vaut pour toute probabilit\'e d'erreur.
Le code Lean4 n'impose que $0 \le \Rel \le 1$.
\end{assumption}

\subsubsection{\'Enonc\'e du th\'eor\`eme original}
\label{subsubsec:enonce_theoreme_original}

\begin{theorem}[Trilemme CRO conditionnel]\label{thm:original_lean4}
Sous les Hypoth\`eses~\ref{ass:fano_axiom},~\ref{ass:pinsker_axiom} et
pour un protocole $\beta$-r\'egulier avec messages uniformes
($H(M)=\lambda$) et $\lambda \ge 64$, on a
\begin{equation}
\Conf(\mathcal{P}) \cdot \Rel(\mathcal{P}) \cdot \Opp(\mathcal{P})
   \le \frac{\Delta}{\lambda} + O(\lambda^{-2}) + \negl(\lambda)
\end{equation}
\end{theorem}

La preuve originale proc\`ede en
quatre \'etapes successives.
\begin{enumerate}
    \item L'in\'egalit\'e de Fano~\cite{Fano61,CoverThomas06} donne
          $I(M;\Sigma|C) \ge \Rel \cdot \lambda - 1$
    \item L'in\'egalit\'e de Pinsker et la
          r\'egularit\'e~\cite{Pinsker1964,Tsybakov2009} donnent
          $\Conf \ge 1 - \sqrt{c_0 \beta I}$
    \item Les hypoth\`eses d'extraction et de
          verdict~\cite{Simon82,Miller56} donnent
          $\Opp \le (I+\Delta)/\gamma$
    \item Le produit est combin\'e sous le r\'egime
          $I \le \kappa/\lambda$ (Hypoth\`ese~\ref{ass:regime}).
\end{enumerate}

\subsubsection{Points critiques et notre approche}
\label{subsubsec:points_critiques}

Notre analyse pr\'eliminaire a identifi\'e trois points qui demandent un
traitement rigoureux.

\paragraph{Statut de $I$ comme variable.}
L'information mutuelle $I(M;\Sigma|C)$ n'est pas une variable libre, puisqu'elle
est d\'etermin\'ee par le protocole. Le r\'egime de pr\'eservation de la
confidentialit\'e $I(\lambda) \le \kappa/\lambda$ (Hypoth\`ese~\ref{ass:regime})
est rendu explicite comme condition du th\'eor\`eme.

\paragraph{Nature asymptotique exacte.}
Nous rempla\c{c}ons le $O(\lambda^{-2})$ informel par une in\'egalit\'e
exacte faisant appara\^itre une constante $K$, une fonction
n\'egligeable $\varepsilon$ et un seuil $N$ explicites. Notre analyse
alg\'ebrique produit la constante $\kappa(\Delta+1) + \kappa^2$,
rendant explicite la d\'ependance en $\kappa$.

\paragraph{Constantes implicites.}
La notation $O(\lambda^{-2})$ dissimule la d\'ependance en
$1/(\alpha H+1)$ et $\beta$. Notre borne explicite
$\gamma^{-1}(\Delta/\lambda + \kappa(\Delta+1)/\lambda^2 + \kappa^2/\lambda^3)$
rend mesurables les param\`etres qui gouvernent la borne~\cite{Grossklags10,Herley09}.

Notre formalisation prouve la borne explicite
$O(\lambda^{-2})$ sous le r\'egime de confidentialit\'e
(Hypoth\`ese~\ref{ass:regime}), enti\`erement v\'erifi\'ee.

\subsubsection{Simplifications et justifications}
\label{subsubsec:simplifications}

\textbf{Simplification~1. Int\'egration des propri\'et\'es informationnelles
comme champs de structure.} Les th\'eor\`emes profonds (Fano, Pinsker) ne sont pas
d\'eclar\'es comme \texttt{axiom} globaux. Ils sont
int\'egr\'es comme propri\'et\'es des instances de protocole, rendant ces
propri\'et\'es falsifiables et v\'erifiables instance par instance.

\textbf{Simplification~2. R\'egime de confidentialit\'e explicite.} Le r\'egime
$I \le \kappa/\lambda$ est pos\'e explicitement comme condition du
th\'eor\`eme, garantissant la confidentialit\'e via le
Lemme~\ref{lem:confidentiality_guarantee}.
\end{sloppypar}

% -----------------------------------------------------------------------
\subsection{D\'efis de formalisation dans Lean4}
\label{subsec:defis_formalisation}
% -----------------------------------------------------------------------
\begin{sloppypar}
Le Trilemme CRO conditionnel pr\'esente plusieurs d\'efis techniques qui
rendent sa v\'erification informelle p\'erilleuse.
Cette sous-section analyse ces d\'efis et les solutions adopt\'ees.

\subsubsection{D\'efi~1~: mod\'eliser les asymptotiques}
\label{subsubsec:defi_asymptotiques}

Le th\'eor\`eme original m\'elange une in\'egalit\'e exacte (valide pour
chaque $\lambda$) avec une notation asymptotique ($O(\lambda^{-2})$,
$\negl(\lambda)$). Pour formaliser cela, nous
introduisons la notion de \textbf{famille de protocoles}, que d\'efinit la
structure suivante.

\begin{leancode}[caption={Famille de protocoles -- structure ProtocolFamily en Lean4}, label={lst:lean4_protocol_family}]
-- Famille de protocoles indexee par le parametre de securite
structure ProtocolFamily where
  params  : InstitutionalParams                -- Parametres institutionnels fixes
  inst    : (lam : ℕ) → ProtocolInstance lam   -- Instance pour chaque lambda
  hparams : ∀ lam, (inst lam).params = params  -- Parametres communs
\end{leancode}

Le th\'eor\`eme devient alors un \'enonc\'e existentiel sur une
constante $K$ et une fonction n\'egligeable $\varepsilon$, valable pour tout
$\lambda \geq 64$, comme l'exprime l'\'enonc\'e ci-dessous.

\begin{leancode}[caption={\'Enonc\'e asymptotique du th\'eor\`eme principal en Lean4}, label={lst:main_theorem_statement}]
-- Theoreme du Trilemme CRO conditionnel
theorem conditional_CRO_trilemma (F : ProtocolFamily) (κ : ℝ) (hκ : 0 ≤ κ)
    (hA : ∀ lam, InstitutionalAssumptions (F.inst lam))
    (hA4 : ∀ lam, (F.inst lam).I ≤ κ / lam) :
    ∃ K : ℝ, ∃ ε : ℕ → ℝ, Negligible ε ∧ ∀ lam : ℕ, 64 ≤ lam →
      (F.inst lam).Conf * (F.inst lam).Rel * (F.inst lam).Opp ≤
        F.params.Δ / lam + K / (lam : ℝ) ^ 2 + ε lam
\end{leancode}

\subsubsection{D\'efi~2~: int\'egrer les r\'esultats profonds}
\label{subsubsec:defi_axiomatiser}

Fano~\cite{Fano61} et Pinsker~\cite{Pinsker1964} sont des th\'eor\`emes
classiques de th\'eorie de l'information. Nous les int\'egrons comme
champs de la structure \texttt{ProtocolInstance}, de sorte que chaque
instance de protocole doit fournir la preuve que ces in\'egalit\'es
s'appliquent \`a ses param\`etres.

\begin{leancode}[caption={Propri\'et\'es informationnelles comme champs de structure}, label={lst:lean4_protocol_instance}]
structure ProtocolInstance (lam : ℕ) where
  params : InstitutionalParams
  Conf   : ℝ   -- confidentialite
  Rel    : ℝ   -- fiabilite
  Opp    : ℝ   -- opposabilite
  I      : ℝ   -- I(M;Σ|C)
  IV     : ℝ   -- I(V;Σ|C)
  HV     : ℝ   -- H(V|C)
  TVmax  : ℝ   -- max TV
  KLmax  : ℝ   -- max D_KL
  hConf  : Conf = 1 - TVmax
  hOpp   : Opp = IV / HV
  hTV    : 0 ≤ TVmax ∧ TVmax ≤ 1
  hRel   : 0 ≤ Rel ∧ Rel ≤ 1
  hI     : 0 ≤ I
  hIV    : 0 ≤ IV
  pinsker : TVmax ^ 2 ≤ c0 * KLmax             -- Lemme 1 (Pinsker)
  fano    : I ≥ Rel * lam - 1                   -- Lemme 2 (Fano), etape 1
\end{leancode}

\subsubsection{D\'efi~3~: r\'egime de confidentialit\'e explicite}
\label{subsubsec:defi_optimisation}

Le r\'egime $I = O(1/\lambda)$ est rendu explicite comme condition du
th\'eor\`eme (Hypoth\`ese~\ref{ass:regime}). La preuve du
Th\'eor\`eme~\ref{thm:trilemme_cro_conditionnel}
utilise directement cette hypoth\`ese sans n\'ecessiter d'optimisation
analytique sur $[0, \lambda]$.

\subsubsection{D\'efi~4~: prouver la n\'egligeabilit\'e}
\label{subsubsec:defi_negligeabilite}

La d\'efinition standard en
cryptographie, \`a savoir
\begin{equation}
f \text{ est n\'egligeable} \iff \forall c>0,\; \exists N,\; \forall n \ge N,\;
|f(n)| \le n^{-c}
\end{equation}
se formalise, avec le filtre \texttt{atTop} de Mathlib (\og pour tout $n$
suffisamment grand\fg~\cite{Mathlib,MathlibAnalysis}) de la mani\`ere suivante.

\begin{leancode}[caption={D\'efinition de n\'egligeabilit\'e en Lean4}, label={lst:lean4_negligible}]
-- Une fonction est negligeable si elle decroit plus vite que tout polynome
def Negligible (f : ℕ → ℝ) : Prop :=
  ∀ c : ℝ, 0 < c → ∀ᶠ n : ℕ in atTop, |f n| ≤ (n : ℝ) ^ (-c)
\end{leancode}

Prouver que $\varepsilon_0(\lambda) = (1/2)^\lambda$ est n\'egligeable
n\'ecessite de montrer que $2^n$ domine $n^c$. La preuve passe par la limite
$n^{\lceil c\rceil}/2^n \to 0$ (lemme \texttt{tendsto\_pow\_const\_div\_const\_pow\_of\_one\_lt}
de Mathlib) et un argument de filtre~; elle ne fournit pas de seuil explicite.
\end{sloppypar}

% -----------------------------------------------------------------------
\subsection{Mod\`ele formel et axiomatisation}
\label{subsec:modele_formel_lean4}
% -----------------------------------------------------------------------
\begin{sloppypar}
\subsubsection{Param\`etres institutionnels}
\label{subsubsec:params_institutionnels_lean4}

Les param\`etres institutionnels sont des constantes fixes pour une
institution donn\'ee, ind\'ependantes de $\lambda$.

\begin{leancode}[caption={Structure InstitutionalParams -- param\`etres institutionnels en Lean4}, label={lst:lean4_institutional_params}]
-- Parametres institutionnels : capacites cognitives et procedurales
structure InstitutionalParams where
  H : ℝ      -- H(C) : entropie contextuelle (bits)
  Δ : ℝ      -- Capacite d'extraction institutionnelle (bits)
  α : ℝ      -- Granularite de verdict (α > 0)
  β : ℝ      -- Constante de regularite du protocole (β ≥ 1)
  hH_nonneg : 0 ≤ H
  hΔ_nonneg : 0 ≤ Δ
  hα_pos    : 0 < α
  hβ_ge_one : 1 ≤ β
\end{leancode}

Cette structure capture les trois familles canoniques~\cite{Adida08,Groth16,Johnson01},
instanci\'ees ci-dessous.

\begin{leancode}[caption={Instances canoniques -- Helios, zk-SNARKs, audit traditionnel}, label={lst:lean4_instances_canoniques}]
noncomputable def HeliosParams : InstitutionalParams :=
  { H := 50, Δ := 56, α := 1.2, β := 1.4, ... }

noncomputable def zkSNARKParams : InstitutionalParams :=
  { H := 70, Δ := 73, α := 1.5, β := 1.5, ... }

noncomputable def TraditionalAuditParams : InstitutionalParams :=
  { H := 20, Δ := 35, α := 1.0, β := 1.0, ... }
\end{leancode}

\subsubsection{Instance de protocole}
\label{subsubsec:instance_protocole_lean4}

\begin{leancode}[caption={Structure ProtocolInstance pour un $\lambda$ donn\'e}, label={lst:lean4_protocol_instance_complete}]
structure ProtocolInstance (lam : ℕ) where
  params : InstitutionalParams
  Conf   : ℝ   -- confidentialite
  Rel    : ℝ   -- fiabilite
  Opp    : ℝ   -- opposabilite
  I      : ℝ   -- I(M;Σ|C)
  IV     : ℝ   -- I(V;Σ|C)
  HV     : ℝ   -- H(V|C)
  TVmax  : ℝ   -- max TV
  KLmax  : ℝ   -- max D_KL
  hConf  : Conf = 1 - TVmax
  hOpp   : Opp = IV / HV
  hTV    : 0 ≤ TVmax ∧ TVmax ≤ 1
  hRel   : 0 ≤ Rel ∧ Rel ≤ 1
  hI     : 0 ≤ I
  hIV    : 0 ≤ IV
  pinsker : TVmax ^ 2 ≤ c0 * KLmax
  fano    : I ≥ Rel * lam - 1
\end{leancode}

\subsubsection{Hypoth\`eses institutionnelles}
\label{subsubsec:axiomes_info_lean4}

\begin{leancode}[caption={Hypoth\`eses institutionnelles en Lean4}, label={lst:lean4_axiomes_extraction}]
structure InstitutionalAssumptions {lam : ℕ} (P : ProtocolInstance lam) : Prop where
  /-- Hypothese 1 (extraction bornee) : I(V;Σ|C) ≤ I(M;Σ|C) + Δ. -/
  A1 : P.IV ≤ P.I + P.params.Δ
  /-- Hypothese 2' (verdict efficace) : H(V|C) ≥ γ = αH(C) − 1 … -/
  A2 : gamma P.params ≤ P.HV
  /-- … avec γ ≥ 1. -/
  A2γ : 1 ≤ gamma P.params
  /-- Hypothese 3 (β-regularite) : D_KL ≤ β·I(M;Σ|C). -/
  A3 : P.KLmax ≤ P.params.β * P.I
\end{leancode}
\end{sloppypar}

% -----------------------------------------------------------------------
\subsection{Preuve du c\oe ur alg\'ebrique}
\label{subsec:preuve_coeur_algebrique}
% -----------------------------------------------------------------------
\begin{sloppypar}
\subsubsection{La constante de Pinsker}
\label{subsubsec:constante_pinsker}

\begin{leancode}[caption={Constante $c_0 = 1/(2\ln 2)$ -- d\'efinition et propri\'et\'es}, label={lst:lean4_c0}]
noncomputable def c0 : ℝ := 1 / (2 * Real.log 2)

lemma c0_pos : 0 < c0 := by
  unfold c0
  have h : 0 < Real.log 2 := Real.log_pos (by norm_num)
  positivity

lemma c0_gt_half : 1 / 2 < c0 := by
  unfold c0
  have h1 := Real.log_two_lt_d9
  have h2 : 0 < Real.log 2 := Real.log_pos (by norm_num)
  rw [lt_div_iff₀ (by positivity)]
  nlinarith
\end{leancode}

\subsubsection{N\'egligeabilit\'e de $\varepsilon_0$}
\label{subsubsec:negligeabilite_eps0}

\begin{lemma}[N\'egligeabilit\'e de $\varepsilon_0$~\cite{Goldreich01,KatzLindell14}]\label{lem:negligeabilite_eps0}
La fonction $\varepsilon_0(\lambda) = (1/2)^\lambda$ est
n\'egligeable.
\end{lemma}

\begin{leancode}[caption={Preuve de n\'egligeabilit\'e en Lean4}, label={lst:lean4_eps0}]
theorem negligible_eps0 : Negligible ε₀ := by
  intro c hc
  have hlim : Tendsto (fun n : ℕ => (n : ℝ) ^ (⌈c⌉₊) / (2 : ℝ) ^ n) atTop (nhds 0) :=
    tendsto_pow_const_div_const_pow_of_one_lt ⌈c⌉₊ (by norm_num)
  have hev : ∀ᶠ n : ℕ in atTop, (n : ℝ) ^ (⌈c⌉₊) / (2 : ℝ) ^ n < 1 :=
    hlim.eventually (gt_mem_nhds one_pos)
  filter_upwards [hev, eventually_ge_atTop 1] with n hn hn1
  -- (suite de la preuve : annexe B, listing complet)
\end{leancode}

\subsubsection{Lemme~4~: garantie de confidentialit\'e}
\label{subsubsec:lemme4_lean4}

\begin{leancode}[caption={Lemme~4 -- garantie de confidentialit\'e en Lean4}, label={lst:lean4_lemma4}]
theorem lemma4_confidentiality_guarantee {lam : ℕ} (P : ProtocolInstance lam)
    (hA3 : P.KLmax ≤ P.params.β * P.I) :
    1 - Real.sqrt (c0 * P.params.β * P.I) ≤ P.Conf := by
  have hTV0 := P.hTV.1
  have h1 : P.TVmax ≤ Real.sqrt (c0 * P.KLmax) := by
    have := Real.abs_le_sqrt P.pinsker
    rwa [abs_of_nonneg hTV0] at this
  have h2 : Real.sqrt (c0 * P.KLmax) ≤ Real.sqrt (c0 * P.params.β * P.I) := by
    apply Real.sqrt_le_sqrt
    rw [mul_assoc]
    exact mul_le_mul_of_nonneg_left hA3 c0_pos.le
  rw [P.hConf]
  linarith
\end{leancode}
\end{sloppypar}

% -----------------------------------------------------------------------
\subsection{Th\'eor\`eme principal}
\label{subsec:theoreme_principal_lean4}
% -----------------------------------------------------------------------

\subsubsection{\'Enonc\'e formel}
\label{subsubsec:enonce_formel_lean4}

\begin{sloppypar}
\begin{theorem}[Trilemme CRO conditionnel~: version v\'erifi\'ee]
\label{thm:conditional_CRO_trilemma}
Soit $F$ une famille de protocoles satisfaisant les conditions suivantes.
\begin{itemize}
    \item les Hypoth\`eses~\ref{ass:extraction_bornee}, \ref{ass:verdicts_efficace}
          et~\ref{ass:regularite} pour chaque instance (structure
          \texttt{InstitutionalAssumptions}, avec $\gamma = \alpha H(C) - 1 \geq 1$)~;
    \item $I_F(\lambda) \leq \kappa/\lambda$ pour un $\kappa \geq 0$
          ind\'ependant de $\lambda$ (r\'egime de pr\'eservation de la
          confidentialit\'e, Hypoth\`ese~\ref{ass:regime})~;
    \item des param\`etres institutionnels $(H(C),\Delta,\alpha,\beta)$ communs \`a
          toute la famille (structure \texttt{ProtocolFamily}).
\end{itemize}
Alors, l'in\'egalit\'e de Fano \'etant appliqu\'ee sous la forme
$I(M;\Sigma|C) \geq \Rel\cdot\lambda - 1$, valable pour tout $\Rel$, il existe une
constante $K$ et une fonction n\'egligeable $\varepsilon$ telles que, pour tout
$\lambda \geq 64$, on ait
\begin{equation}
\Conf_F(\lambda) \cdot \Rel_F(\lambda) \cdot \Opp_F(\lambda) \leq \frac{\Delta}{\lambda} + \frac{K}{\lambda^2} + \varepsilon(\lambda)
\label{eq:main_theorem_lean4}
\end{equation}
Instance par instance (\texttt{trilemma\_explicit}), on obtient la forme avec le facteur $1/\gamma$
\begin{equation}
\Conf_F(\lambda) \cdot \Rel_F(\lambda) \cdot \Opp_F(\lambda) \leq \frac{1}{\gamma}\left(\frac{\Delta}{\lambda} + \frac{\kappa(\Delta+1)}{\lambda^2} + \frac{\kappa^2}{\lambda^3}\right)
\label{eq:main_theorem_lean4_explicit}
\end{equation}
\end{theorem}
\end{sloppypar}

\subsubsection{T\'emoins de coh\'erence et de port\'ee}
\label{subsubsec:temoins_lean4}

\begin{sloppypar}
La formalisation Lean4 produit deux t\'emoins qui d\'emontrent la
non-vacuit\'e et la n\'ecessit\'e des hypoth\`eses du th\'eor\`eme.

\begin{theorem}[Coh\'erence du mod\`ele abstrait]
\label{thm:witness_consistency}
Il existe une instance du mod\`ele abstrait v\'erifiant simultan\'ement toutes
les hypoth\`eses, avec $\Conf = \Opp = 1$, $\Rel = 1/\lambda$, $I = 0$,
$\Delta = \gamma = 1$. Pour cette instance, on a
\begin{equation}
\Conf \cdot \Rel \cdot \Opp = \frac{1}{\lambda} = \frac{\Delta}{\gamma \lambda}
\end{equation}
Cette valeur \'egale exactement la borne du
Th\'eor\`eme~\ref{thm:conditional_CRO_trilemma}. Cela d\'emontre que les
hypoth\`eses sont non vides et que la borne est atteinte \emph{dans le
mod\`ele abstrait}.
\end{theorem}

\begin{remark}[Port\'ee du t\'emoin de coh\'erence]\label{rem:portee_temoin}
Une \texttt{ProtocolInstance} est un point du mod\`ele abstrait, c'est-\`a-dire un ensemble de r\'eels
qui satisfont les contraintes de Pinsker, de Fano (sous la forme rel\^ach\'ee
$I \geq \Rel\cdot\lambda - 1$) et les Hypoth\`eses~1 \`a~4. Le t\'emoin \'etablit la
non-vacuit\'e de ce mod\`ele. L'existence d'une loi de probabilit\'e qui le
r\'ealise est une autre question. Pour une loi r\'eelle avec $M$ uniforme sur $\lambda$ bits, $I = 0$
impose $\Relrec = 2^{-\lambda}$ au lieu de $1/\lambda$. Le t\'emoin atteint donc la borne
rel\^ach\'ee~; il ne montre pas qu'un protocole r\'eel l'atteint. Il a $\Delta = \gamma = 1$
et r\'ealise l'\'egalit\'e dans la forme rel\^ach\'ee de Fano. Il illustre ainsi le r\'egime
$\Delta \geq \gamma$ de la Remarque~\ref{rem:activite_delta}. Le t\'emoin de port\'ee
(Th\'eor\`eme~\ref{thm:scope_necessity}) est, lui, r\'ealis\'e par une loi explicite
(Remarque~\ref{rem:scope_trilemme_classique}).
\end{remark}

\begin{leancode}[caption={T\'emoin de coh\'erence~: instance atteignant la borne}, label={lst:lean4_witness_instance}]
theorem assumptions_consistent_and_bound_attained (lam : ℕ) (hlam : 1 ≤ lam) :
    InstitutionalAssumptions (witnessInstance lam hlam) ∧
    (witnessInstance lam hlam).I ≤ 0 / lam ∧
    (witnessInstance lam hlam).Conf * (witnessInstance lam hlam).Rel *
      (witnessInstance lam hlam).Opp =
      (witnessInstance lam hlam).params.Δ / (gamma (witnessInstance lam hlam).params * lam) := by
  refine ⟨⟨?_, ?_, ?_, ?_⟩, ?_, ?_⟩ <;>
    simp [witnessInstance, WitnessParams, gamma] <;> norm_num
\end{leancode}

\begin{theorem}[Port\'ee~: n\'ecessit\'e du r\'egime]
\label{thm:scope_necessity}
Il existe une instance de protocole v\'erifiant les Hypoth\`eses~1, 2$'$
et 3 avec $\Delta = 0$ mais hors du r\'egime de l'Hypoth\`ese~\ref{ass:regime}
($I = \lambda/2 - 1$). Pour cette instance, on a
\begin{equation}
\Conf \cdot \Rel \cdot \Opp = \frac{1}{4}
\end{equation}
pour tout $\lambda \geq 64$, ce produit restant constant. Cela d\'emontre que le r\'egime
$I \leq \kappa/\lambda$ est n\'ecessaire \`a la borne.
\end{theorem}

\begin{leancode}[caption={T\'emoin de port\'ee~: instance hors du r\'egime}, label={lst:lean4_scope_instance}]
theorem regime_assumption_necessary (lam : ℕ) (hlam : 64 ≤ lam) :
    InstitutionalAssumptions (scopeInstance lam hlam) ∧
    (scopeInstance lam hlam).params.Δ = 0 ∧
    (scopeInstance lam hlam).Conf * (scopeInstance lam hlam).Rel *
      (scopeInstance lam hlam).Opp = 1 / 4 ∧
    (scopeInstance lam hlam).I = lam / 2 - 1 := by
  have hl : (64 : ℝ) ≤ lam := by exact_mod_cast hlam
  refine ⟨⟨?_, ?_, ?_, ?_⟩, ?_, ?_, ?_⟩ <;>
    simp [scopeInstance, ScopeParams, gamma] <;> linarith
\end{leancode}

\begin{remark}[Statut de v\'erification]
\label{rem:statut_verification_lean4}
Le noyau alg\'ebrique de cette preuve est certifi\'e m\'ecaniquement en
Lean4. Les formulations alternatives issues de la certification sont
pr\'esent\'ees \`a la sous-section~\ref{subsec:lean4_erreurs}.
\end{remark}
\end{sloppypar}

% -----------------------------------------------------------------------
\subsection{Extensions~: verdict d\'eterministe et forme raffin\'ee}
\label{subsec:lean4_extensions}
% -----------------------------------------------------------------------
\begin{sloppypar}
Quatre \'enonc\'es compl\`etent la formalisation (version~2.2 du fichier,
listing~\ref{lst:lean4_extensions})~; ils certifient les deux faits
structurels \'etablis au chapitre~\ref{chap:cadre_formel}.
\begin{itemize}
\item \texttt{opp\_eq\_one\_of\_deterministic} \'etablit que, si $I(V;\Sigma|C) = H(V|C) > 0$
      (verdict d\'eterministe), alors $\Opp = 1$
      (Remarque~\ref{rem:verdict_deterministe})~;
\item \texttt{deterministic\_forces\_capacity} montre que, dans ce cas, les
      Hypoth\`eses~1 et~2$'$ imposent $\gamma - I(M;\Sigma|C) \leq \Delta$~;
\item \texttt{trilemma\_refined} certifie la forme raffin\'ee~\eqref{eq:trilemme_cro_raffine},
      sous la condition $I(V;\Sigma|C) \leq H(V|C)$, toujours vraie~;
\item \texttt{refined\_bound\_is\_fano\_when\_delta\_ge\_gamma} \'etablit que, si
      $\Delta \geq \gamma$, le facteur institutionnel vaut~$1$ et la borne se
      r\'eduit \`a celle de Fano (Remarque~\ref{rem:activite_delta}).
\end{itemize}
Chacun de ces \'enonc\'es ne d\'epend que des axiomes standard \texttt{propext},
\texttt{Classical.choice} et \texttt{Quot.sound}.

\begin{leancode}[caption={Extensions de la version 2.2}, label={lst:lean4_extensions}]
/-- **T1.** Verdict déterministe : I(V;Σ|C) = H(V|C) ⇒ Opp = 1. -/
theorem opp_eq_one_of_deterministic {lam : ℕ} (P : ProtocolInstance lam)
    (hdet : P.IV = P.HV) (hHV : 0 < P.HV) : P.Opp = 1 := by
  rw [P.hOpp, hdet]; exact div_self hHV.ne'

/-- **T1 bis.** Verdict déterministe + Hyp. 1 et 2′ ⇒ γ − I(M;Σ|C) ≤ Δ :
    la contrainte institutionnelle ne peut pas être active. -/
theorem deterministic_forces_capacity {lam : ℕ} (P : ProtocolInstance lam)
    (hA : InstitutionalAssumptions P) (hdet : P.IV = P.HV) :
    gamma P.params - P.I ≤ P.params.Δ := by
  have := hA.A1; have := hA.A2; linarith

/-- **T2 (forme raffinée).** Avec I(V;Σ|C) ≤ H(V|C) (toujours vrai) :
    Conf·Rel·Opp ≤ ((κ/λ+1)/λ) · min 1 ((κ/λ+Δ)/γ). -/
theorem trilemma_refined {lam : ℕ} (P : ProtocolInstance lam)
    (hA : InstitutionalAssumptions P) (κ : ℝ) (hκ : 0 ≤ κ)
    (hlam : 1 ≤ lam) (hA4 : P.I ≤ κ / lam) (hIVHV : P.IV ≤ P.HV) :
    P.Conf * P.Rel * P.Opp ≤
      ((κ / lam + 1) / lam) * min 1 ((κ / lam + P.params.Δ) / gamma P.params) := by
  have hl0 : (0 : ℝ) < lam := by exact_mod_cast hlam
  have hγ0 : 0 < gamma P.params := by linarith [hA.A2γ]
  have hHV : 0 < P.HV := by linarith [hA.A2, hA.A2γ]
  have hκl : 0 ≤ κ / lam := div_nonneg hκ hl0.le
  have hC1 : P.Conf ≤ 1 := by rw [P.hConf]; linarith [P.hTV.1]
  have hC0 : 0 ≤ P.Conf := by rw [P.hConf]; linarith [P.hTV.2]
  have hR0 := P.hRel.1
  have hR : P.Rel ≤ (κ / lam + 1) / lam := by
    rw [le_div_iff₀ hl0]; linarith [P.fano]
  have hO0 : 0 ≤ P.Opp := by rw [P.hOpp]; exact div_nonneg P.hIV hHV.le
  have hO1 : P.Opp ≤ 1 := by
    rw [P.hOpp, div_le_one hHV]; exact hIVHV
  have hO2 : P.Opp ≤ (κ / lam + P.params.Δ) / gamma P.params := by
    rw [P.hOpp, div_le_div_iff₀ hHV hγ0]
    have e1 := mul_le_mul_of_nonneg_right hA.A1 hγ0.le
    have e2 := mul_le_mul_of_nonneg_left hA.A2
      (by linarith [P.params.hΔ_nonneg] : (0 : ℝ) ≤ κ / lam + P.params.Δ)
    have e3 : P.I * gamma P.params ≤ κ / lam * gamma P.params :=
      mul_le_mul_of_nonneg_right hA4 hγ0.le
    nlinarith
  have hOm : P.Opp ≤ min 1 ((κ / lam + P.params.Δ) / gamma P.params) := le_min hO1 hO2
  calc P.Conf * P.Rel * P.Opp ≤ 1 * ((κ / lam + 1) / lam) * min 1 ((κ / lam + P.params.Δ) / gamma P.params) :=
        mul_le_mul (mul_le_mul hC1 hR hR0 zero_le_one) hOm hO0
          (by rw [one_mul]; exact div_nonneg (by linarith) hl0.le)
    _ = _ := by ring

/-- **T2 bis.** Si Δ ≥ γ, la borne raffinée coïncide avec la seule borne de Fano. -/
theorem refined_bound_is_fano_when_delta_ge_gamma {lam : ℕ} (P : ProtocolInstance lam)
    (hA : InstitutionalAssumptions P) (κ : ℝ) (hκ : 0 ≤ κ) (hlam : 1 ≤ lam)
    (hΔγ : gamma P.params ≤ P.params.Δ) :
    min 1 ((κ / lam + P.params.Δ) / gamma P.params) = 1 := by
  have hl0 : (0 : ℝ) < lam := by exact_mod_cast hlam
  have hγ0 : 0 < gamma P.params := by linarith [hA.A2γ]
  have hκl : 0 ≤ κ / lam := div_nonneg hκ hl0.le
  apply min_eq_left
  rw [le_div_iff₀ hγ0]; linarith
\end{leancode}
\end{sloppypar}

% -----------------------------------------------------------------------
\subsection{Ce que la formalisation a r\'ev\'el\'e}
\label{subsec:lean4_erreurs}
% -----------------------------------------------------------------------
\begin{sloppypar}
La formalisation Lean4 du noyau alg\'ebrique du Trilemme CRO a produit,
au-del\`a de la certification, des formulations alternatives de plusieurs
\'etapes des preuves publi\'ees~\cite{MINKA2026CRO,MINKA2026PQCRO}. Ces
formulations, plus pr\'ecises que les versions initiales, \'eclairent la
structure logique du r\'esultat et d\'elimitent son domaine de validit\'e.
Elles sont document\'ees ci-dessous, car elles constituent une contribution
m\'ethodologique \`a part enti\`ere.

\paragraph{Sens de l'in\'egalit\'e du Lemme~4.}
Le Lemme~4 publi\'e formulait sur $\Conf$ la borne \emph{sup\'erieure}
$\Conf \leq 1 - \sqrt{c_0 \beta I / \lambda}$. La formalisation a conduit
\`a une formulation alternative, dans laquelle l'\'enonc\'e prend la forme
d'une borne \emph{inf\'erieure} $\Conf \geq 1 - \sqrt{c_0 \beta I}$. Cette
formulation pr\'ecise la port\'ee de l'argument de Pinsker, qui ne requiert
pas de normalisation par $\lambda$. Le
Lemme~\ref{lem:confidentiality_guarantee} retient cette formulation
alternative.

\paragraph{Statut de l'Hypoth\`ese~2.}
L'Hypoth\`ese~2 publi\'ee \'enon\c{c}ait $|\mathcal{Y}_\lambda| \leq
2^{\alpha H(C)}$ (borne sup\'erieure) et en d\'eduisait
$H(V|C) \geq \alpha H(C) - 1$ (borne inf\'erieure). La formalisation a
montr\'e que ces deux \'enonc\'es sont logiquement ind\'ependants. La borne
inf\'erieure exprime une propri\'et\'e d'\emph{efficacit\'e} du verdict, que
la borne sup\'erieure ne suffit pas \`a \'etablir. La formulation alternative
retenue s\'epare explicitement les deux \'enonc\'es
(Hypoth\`ese~\ref{ass:verdicts_efficace}).

\paragraph{R\'esolution de l'optimisation.}
L'\'Etape~6 de la preuve publi\'ee r\'esolvait $f'(x) = 0$ avec
$f'(x) = 2x - \frac{5}{2}ax^{3/2}$, en retenant $\sqrt{x^*} =
\frac{5a}{4}$. La formalisation a conduit \`a une r\'esolution
alternative, $\sqrt{x^*} = \frac{4}{5a}$. Elle a montr\'e que ce point
est un \emph{maximiseur} de $f$. Le r\'egime $I = O(1/\lambda)$ ne
s'obtient donc pas par cette optimisation. La formulation alternative
pose ce r\'egime comme hypoth\`ese explicite (Hypoth\`ese~\ref{ass:regime}).

\paragraph{N\'ecessit\'e du r\'egime~: un cas limite.}
La formalisation a mis en \'evidence un cas limite qui \'etablit la
n\'ecessit\'e \emph{logique} du r\'egime de confidentialit\'e
(Hypoth\`ese~\ref{ass:regime}). Sans ce r\'egime, la borne du Trilemme
cesse d'\^etre valide, pour une raison structurelle que le cas ci-dessous
rend explicite.

Consid\'erons le protocole suivant, en encodage classique diagonal
($\Sigma$ prend ses valeurs dans un alphabet fini $\mathcal{S}$). Le
protocole r\'ev\`ele $M$ avec probabilit\'e $1/2$ et produit sinon une
preuve uniform\'ement al\'eatoire dans $\{0,1\}^\lambda$, ce qui s'\'ecrit
\begin{equation}
\Sigma =
\begin{cases}
M, & \text{avec probabilit\'e } 1/2, \\
U \xleftarrow{\$} \{0,1\}^\lambda, & \text{avec probabilit\'e } 1/2.
\end{cases}
\label{eq:counterexample_protocol}
\end{equation}
Soit $V$ le premier bit de $\Sigma$. Nous supposons de plus $M \perp C$.

Pour ce protocole, les trois m\'etriques valent
\begin{align}
\Conf &= 1/2, \\
\Rel &= 1/2 + 2^{-\lambda-1}, \\
\Opp &= 1,
\end{align}
et $I(M;\Sigma|C) = \lambda/2 - 1 + O(\lambda 2^{-\lambda})$. Le
protocole est $\beta$-r\'egulier avec $\beta \leq 1{,}04$ pour tout
$\lambda \geq 64$. Il satisfait donc les
Hypoth\`eses~\ref{ass:extraction_bornee}--\ref{ass:regularite}. Il
\emph{s'\'ecarte} cependant de l'Hypoth\`ese~\ref{ass:regime}, car
$I(M;\Sigma|C) = \Theta(\lambda)$ alors que le r\'egime exige $O(1/\lambda)$.

Pour ce protocole, le produit CRO vaut
\begin{equation}
\Conf \cdot \Rel \cdot \Opp \;\approx\; \frac{1}{2} \cdot \frac{1}{2} \cdot 1 \;=\; \frac{1}{4},
\end{equation}
ind\'ependant de $\lambda$. La borne du Trilemme
$\Conf \cdot \Rel \cdot \Opp \leq \Delta/\lambda + O(\lambda^{-2})$ ne
tient donc pas uniform\'ement pour $I = \Theta(\lambda)$, puisqu'elle
requiert le r\'egime de confidentialit\'e.

\textbf{Cons\'equence.} Ce cas limite \'etablit que
l'Hypoth\`ese~\ref{ass:regime} est une \emph{condition n\'ecessaire} de la
borne. Il justifie r\'etrospectivement la formulation alternative du
Th\'eor\`eme~\ref{thm:trilemme_cro_conditionnel}, dans laquelle cette
hypoth\`ese est explicite. Il a aussi motiv\'e la distinction formelle, dans
l'annexe~\ref{chap:annexe_lean4}, entre le \textbf{t\'emoin de coh\'erence}
(qui atteint la borne) et le \textbf{t\'emoin de port\'ee} (qui s'\'ecarte du
r\'egime et montre que la borne ne tient pas sans lui).

\paragraph{Propagation \`a l'extension quantique.}
Les formulations alternatives de la preuve classique ont conduit \`a
r\'eexaminer l'extension post-quantique du Trilemme~\cite{MINKA2026PQCRO},
qui h\'eritait de la m\^eme structure. Le Lemme~6 et la borne sup\'erieure
du Lemme~7 ont fait l'objet de formulations plus pr\'ecises. Pour le
Lemme~6, la formalisation a mis en \'evidence un cas limite form\'e de deux
\'etats orthogonaux. Les Hypoth\`eses~4 et~5 sont d\'esormais pos\'ees
explicitement au lieu d'\^etre d\'eduites. Les formulations correspondantes
sont document\'ees aux
sections~\ref{subsec:enonce_preuve_pq} et~\ref{subsec:hypotheses_quantiques}.

\paragraph{Comment la formalisation a produit ces formulations.}
Ces formulations alternatives sont issues de la formalisation. La relecture
des articles ne les avait pas fait appara\^itre. Elles sont n\'ees \`a deux
moments pr\'ecis. D'abord, l'\'enonc\'e publi\'e du Lemme~4 n'a pas pu \^etre
d\'emontr\'e \`a partir des hypoth\`eses disponibles, ce qui a conduit \`a
r\'eexaminer la preuve et \`a pr\'eciser le sens de l'in\'egalit\'e. Un \'echec
de preuve formelle ne tranche pas \`a lui seul~; l'examen qui l'a suivi a
\'etabli la formulation retenue. Ensuite, la formalisation du noyau
alg\'ebrique a montr\'e que l'optimisation publi\'ee ne pouvait pas \^etre
reproduite. Le cas limite ci-dessus a \'emerg\'e de l'effort pour pr\'eserver
la coh\'erence de la preuve formelle. Il est alors devenu n\'ecessaire
d'inscrire explicitement le r\'egime de confidentialit\'e dans l'\'enonc\'e du
th\'eor\`eme.

\paragraph{Statut des formulations alternatives.}
Les formulations alternatives des preuves classique et post-quantique sont
celles pr\'esent\'ees aux sections~\ref{sec:cro_classique}
et~\ref{sec:cro_pq} de cette th\`ese. Les versions publi\'ees initiales
restent accessibles pour r\'ef\'erence~\cite{MINKA2026CRO,MINKA2026PQCRO}.

\paragraph{La le\c{c}on m\'ethodologique.}
Les formulations alternatives et le cas limite qui \'etablit la n\'ecessit\'e
du r\'egime de confidentialit\'e conduisent \`a une le\c{c}on
m\'ethodologique sur le r\^ole de la formalisation dans le processus de
publication scientifique, d\'evelopp\'ee \`a la
section~\ref{subsec:verification_corrige_lecon}.
\end{sloppypar}

\begin{boiteArbitre}
L'arbitre referme le dossier. La machine a produit des formulations
alternatives que les auteurs ont retenues. La le\c{c}on qu'il en tire
d\'epasse la substitution d'une formule \`a une autre, car ni la relecture
humaine ni la machine ne suffisent seules. La machine v\'erifie les
d\'eductions tandis que la relecture v\'erifie que les d\'efinitions disent ce
qu'on croit. L'arbitre comprend que la formalisation est une n\'ecessit\'e
\'epist\'emologique pour tout r\'esultat interdisciplinaire dont aucune des
disciplines constituantes ne contr\^ole l'int\'egralit\'e.
\end{boiteArbitre}


% -----------------------------------------------------------------------
\subsection{Application aux familles de protocoles canoniques}
\label{subsec:application_familles_canoniques}
% -----------------------------------------------------------------------
\begin{sloppypar}
\subsubsection{R\'egimes de capacit\'e institutionnelle}
\label{subsubsec:regimes_capacite_lean4}

La formalisation contient aussi une classification indicative des
capacit\'es institutionnelles par le rapport $\Delta/\lambda$. Elle est distincte
des r\'egimes du th\'eor\`eme (seuil $H^\sharp$, Remarque~\ref{rem:activite_delta})
et des r\'egimes de construction SOTP (section~\ref{sec:sotp_construction})~; elle
sert seulement \`a illustrer l'usage de d\'efinitions calculables dans Lean. Les
param\`etres des familles canoniques cod\'ees en Lean sont ceux de la
calibration illustrative, \`a l'exception de $H(C)$, qui n'intervient pas dans
cette classification.

\begin{leancode}[caption={Classification des r\'egimes institutionnels en Lean4}, label={lst:lean4_capacity_regime}]
inductive InstitutionalRegime where
  | LowCapacity    -- Δ < 0.1 × λ : capacite faible
  | MidCapacity    -- 0.1 × λ ≤ Δ ≤ 0.5 × λ : capacite moyenne
  | HighCapacity   -- Δ > 0.5 × λ : capacite elevee

noncomputable def capacity_regime (Δ lam : ℝ) : InstitutionalRegime :=
  if Δ < 0.1 * lam then .LowCapacity
  else if Δ ≤ 0.5 * lam then .MidCapacity
  else .HighCapacity
\end{leancode}

\begin{theorem}[R\'egime Helios~\cite{Adida08}]\label{thm:regime_helios}
\`A $\lambda = 128$, Helios ($\Delta = 56$) est
dans le r\'egime de capacit\'e moyenne, car $56 \in [12{,}8,\, 64]$.
\end{theorem}

\begin{leancode}[caption={V\'erification du r\'egime Helios en Lean4}, label={lst:lean4_helios_regime}]
theorem Helios_regime_at_128 :
    capacity_regime HeliosParams.Δ 128 = .MidCapacity := by
  show capacity_regime 56 128 = _
  unfold capacity_regime
  norm_num
\end{leancode}

\begin{theorem}[R\'egime zk-SNARKs~\cite{Groth16,BCTV14}]\label{thm:regime_snarks}
\`A $\lambda = 256$, zk-SNARKs ($\Delta = 73$) est
dans le r\'egime de capacit\'e moyenne, car $73 \in [25{,}6,\, 128]$.
\end{theorem}

\begin{leancode}[caption={V\'erification du r\'egime zk-SNARKs en Lean4}, label={lst:lean4_zksnark_regime}]
theorem zkSNARK_regime_at_256 :
    capacity_regime zkSNARKParams.Δ 256 = .MidCapacity := by
  show capacity_regime 73 256 = _
  unfold capacity_regime
  norm_num
\end{leancode}

\begin{theorem}[R\'egime audit traditionnel~\cite{Mason07,Kerr04}]\label{thm:regime_audit}
\`A $\lambda = 128$, l'audit traditionnel ($\Delta = 35$) est dans le r\'egime de capacit\'e moyenne, car
$35 \in [12{,}8,\, 64]$.
\end{theorem}

\begin{leancode}[caption={V\'erification du r\'egime audit traditionnel en Lean4}, label={lst:lean4_audit_regime}]
theorem TraditionalAudit_regime_at_128 :
    capacity_regime TraditionalAuditParams.Δ 128 = .MidCapacity := by
  show capacity_regime 35 128 = _
  unfold capacity_regime
  norm_num
\end{leancode}

Ces v\'erifications sont enti\`erement automatis\'ees dans Lean4. Elles
v\'erifient des calculs sur des valeurs illustratives, sans valider le mod\`ele
de capacit\'e institutionnelle.
\end{sloppypar}

% -----------------------------------------------------------------------
% Bo\^ite finale
% -----------------------------------------------------------------------
\begin{boiteArbitre}
L'arbitre referme le dossier Lean4. Il a achev\'e l'\'etape \textbf{C}.
Cette section lui a appris que la preuve qu'il croyait compl\`ete gagnait
\`a \^etre pr\'ecis\'ee en plusieurs points.

Sept contributions structurent ce qu'il retient. Il y a d'abord le mod\`ele formel
(familles de protocoles index\'ees par $\lambda$, param\`etres institutionnels
fixes), puis l'int\'egration des propri\'et\'es informationnelles comme champs de
structure (Fano, Pinsker) et le noyau alg\'ebrique certifi\'e sans aucune
admission tacite. S'y ajoutent la preuve de n\'egligeabilit\'e, l\`a o\`u l'expos\'e
informel \'ecrivait simplement \og n\'egligeable\fg, le r\'egime de confidentialit\'e
rendu explicite comme condition du th\'eor\`eme ainsi que les constantes explicites
($\kappa(\Delta+1)$ et $\kappa^2$ au lieu d'un $O(\lambda^{-2})$ non chiffr\'e). La
forme raffin\'ee compl\`ete l'ensemble en lui indiquant exactement quand sa capacit\'e
d'extraction devient la contrainte liante.

L'encha\^inement des faits le marque particuli\`erement. Trois formulations
initiales (le sens d'une in\'egalit\'e, une borne inf\'erieure tir\'ee d'une borne
sup\'erieure, la r\'esolution d'une optimisation) ont \'et\'e pr\'ecis\'ees, ce que
la relecture des articles n'avait pas permis. La formalisation a oblig\'e \`a
justifier chaque \'etape. Le cas limite qui \'etablit la n\'ecessit\'e du
r\'egime est d'ailleurs n\'e de cette exigence.

La borne $\Delta/(\gamma\lambda) + O(\lambda^{-2})$ est maintenant prouv\'ee dans le
cadre du r\'egime, avec des termes correctifs chiffr\'es. Cette section en tire sa
port\'ee \'epist\'emologique, puisqu'elle montre que la rigueur se v\'erifie au lieu
de se supposer.
\end{boiteArbitre}